%% ==========================================================================
%%  FINAL REPORT --- SECTION 3, BACKGROUND
%%
%%  WHY THIS IS A NEW FILE AND NOT %% ==========================================================================
%%  SECTION 4 --- CRYPTOGRAPHIC BUILDING BLOCKS
%%  Every claim traces to report/notes/evidence.md.
%% ==========================================================================
\section{Cryptographic Building Blocks}
\label{sec:prim}

This section introduces the primitives the systems in Sections~\ref{sec:training}
and~\ref{sec:retrieval} are assembled from. The treatment is deliberately shallow on
constructions and deliberate about \emph{costs}, because it is the cost profile of each primitive,
and in particular whether it costs communication, computation, or rounds, that determines
which systems are built from which pieces.

\subsection{Secret sharing}
\label{sec:prim:sharing}

Secret sharing splits a value among parties so that authorised subsets can reconstruct it and
unauthorised subsets learn nothing~\cite{shamir1979}. Two variants recur throughout this survey.

In \emph{additive} sharing over a ring $\ring = \Ztwo{\bits}$, a value $x$ is split as
$x = x_0 + x_1 \bmod 2^{\bits}$, with one summand given to each of two parties. Addition of shared
values, and multiplication by a public scalar, are then local operations requiring no
communication. Multiplication of two shared values is not, and is classically handled with
precomputed correlated randomness in the form of multiplication triples~\cite{beaver1991}.

In \emph{replicated} $2$-out-of-$3$ sharing, $x = x_0 + x_1 + x_2$ and party $i$ holds the pair
$(x_i, x_{i+1})$. Any two parties can reconstruct it, while a single party learns nothing. This is the
sharing used by the three-party systems in Section~\ref{sec:training:mpc}, and its attraction is
that a party holding two of the three summands can compute a share of a product locally, so that
degree-two functions become very cheap~\cite{araki2016}.

Several systems additionally assume the \emph{PRF model}, in which each pair of parties holds a
shared pseudorandom-function key and can therefore generate correlated randomness, notably
shares of zero, without communicating at all.

\subsection{Garbled circuits}
\label{sec:prim:gc}

Garbled circuits~\cite{yao1986} allow two parties to evaluate a Boolean circuit on private inputs
in a constant number of rounds. The cost is borne in bandwidth: the garbled circuit must be
transmitted, so communication scales with circuit size. For data-oblivious algorithms whose
circuit size grows with the whole input, as is the case for gradient descent over a rating
matrix, this is the binding constraint, and it is the reason the systems in
Section~\ref{sec:training:gc} are superseded by later secret-sharing approaches.

\subsection{Homomorphic encryption}
\label{sec:prim:he}

Homomorphic encryption permits computation directly on ciphertexts. Partially homomorphic schemes
support one operation. Additively homomorphic schemes such as Paillier's~\cite{paillier1999}
are used in several of the systems below, while fully homomorphic encryption supports
arbitrary circuits. Gentry~\cite{gentry2009} established the latter is possible via
\emph{bootstrapping}: it suffices to construct a scheme that can homomorphically evaluate a
slightly augmented version of its own decryption circuit.

The relevance of homomorphic encryption to this survey is structural rather than technical. It is
the primitive that permits \emph{single-server} constructions, and therefore the one that removes
the non-collusion assumption discussed in Section~\ref{sec:taxonomy:trust}, at a computational
cost that Sections~\ref{sec:prim:pir} and~\ref{sec:retrieval} quantify.

\subsection{Function secret sharing and distributed point functions}
\label{sec:prim:fss}

Function secret sharing (FSS) shares a \emph{function} rather than a value. Following Boyle et
al.~\cite{bgi2016}, an $m$-party FSS scheme splits a function $f\colon \{0,1\}^n \to G$, for an
abelian group $G$, into keys $k_1, \dots, k_m$ such that $f = f_1 + \dots + f_m$ and every strict
subset of the keys hides $f$. A \emph{distributed point function} (DPF) is the special case in
which $f$ is a point function: non-zero at a single input and zero everywhere
else~\cite{gilboa2014}.

The applications are named directly by the authors: FSS schemes are useful for ``privately reading
from or writing to distributed databases while minimizing the amount of communication'',
including ``different flavors of private information retrieval (PIR), as well as a recent
application of DPF for large-scale anonymous messaging''~\cite{bgi2016}. The same work reports
reducing the key size of the earlier PRG-based DPF construction ``roughly by a factor of
4''~\cite{bgi2016}.

The property that matters for everything downstream is \emph{compression}. A DPF key is
logarithmic in the domain size, so a party can hold a short key and expand it into a long
pseudorandom vector that differs from its counterpart at exactly one secret position. This is
what makes the read layer of Section~\ref{sec:prim:pir} and the DORAM of
Section~\ref{sec:prim:oram} practical.

FSS also generalises beyond point functions into a technique for secure computation. Boyle et
al.~\cite{boyle2021} evaluate a gate $g$ using an FSS scheme for the offset family
$g_r(x) = g(x + r)$, giving schemes for ``zero test, integer comparison, ReLU, and spline
functions'' and reporting ``significant savings in online communication and round complexity
compared to alternative techniques based on garbled circuits or secret sharing'', along with
``roughly $4\times$ reduction in key size'' for the distributed comparison function. Subsequent
work builds function libraries on this foundation: LLAMA~\cite{llama2022} operates in the trusted
dealer model, ``where a dealer provides input independent correlated randomness to both
parties'', yielding ``a very lightweight online phase'', while noting that prior FSS-based
two-party work lacked ``support for math functions (e.g., sigmoid, and reciprocal square root)''
and restricted all values to a single bitwidth. Pika~\cite{pika2022} reports protocols for
division, exponentiation, logarithm and tanh with ``constant round complexity (3 for semi-honest,
4 for malicious)'' and communication ``54--121$\times$ lower than prior art'', and, unusually
for this survey, extends them to malicious security in the honest-majority setting.
\grotto{} contributes a structural observation about the DPF tree that lets ``a single DPF'' do
what state-of-the-art approaches ``require many DCFs to do''.

\subsection{Private information retrieval}
\label{sec:prim:pir}

Private information retrieval (PIR) lets a client fetch record $j$ from a server-held database
without the server learning $j$. The literature divides on the trust assumption.

\emph{Multi-server} PIR distributes the database across servers assumed not to
collude~\cite{chor1995}. Hafiz and Henry~\cite{hafiz2019} study this setting, describing ``several
untrusted servers work to obliviously service remote clients' requests for data and yet no pair of
servers colludes'', and report efficiency ``with respect to every cost metric --- download,
upload, computation, and round complexity''. Their family of protocols recovers, as special
cases, the two-server instances of the seminal scheme of Chor et al.\ and the DPF-based
construction of Boyle et al.~\cite{hafiz2019}. This is the scheme \pirsona{} builds on.

\emph{Single-server} PIR removes the non-collusion assumption at a computational
price~\cite{kushilevitz1997}. SimplePIR~\cite{simplepir2023} is presented as ``the fastest
single-server private information retrieval scheme known to date'', with security under
learning-with-errors, reporting ``10 GB/s/core server throughput'' that ``approaches \dots\ the
performance of the fastest two-server private-information-retrieval schemes (which require
non-colluding servers)''. The authors state the cost plainly: ``relatively large communication
costs: to make queries to a 1 GB database, the client must download a 121 MB `hint' about the
database contents'', after which each query needs 242~KB. A variant, DoublePIR, shrinks the hint
to 16~MB at 345~KB per query and 7.4~GB/s/core~\cite{simplepir2023}. Spiral~\cite{spiral2022}
composes two lattice-based encryption schemes and reports ``at least a $4.5\times$ reduction in
query size, $1.5\times$ reduction in response size, and $2\times$ increase in server throughput
compared to previous systems''.

Two further directions matter later. Corrigan-Gibbs and Kogan~\cite{cgk2020} construct
single-server schemes with ``sublinear amortized server time'' by having the client ``first fetch
a small `hint' about the database contents'', paying a linear cost once and answering subsequent
queries in sublinear time. This preprocessing paradigm is what Pacmann (Section~\ref{sec:retrieval:sublinear})
depends on. Angel et al.~\cite{sealpir2018} reduce query size ``up to $274\times$'' and introduce
probabilistic batch codes giving ``up to $40\times$ speedup over processing queries one at a
time'', which matters for any recommender serving many users concurrently.

\subsection{Oblivious RAM and distributed ORAM}
\label{sec:prim:oram}

Oblivious RAM hides \emph{access patterns}, not merely contents~\cite{goldreich1996}. Path
ORAM~\cite{pathoram2013} is the canonical tree-based construction, ``an extremely simple Oblivious
RAM protocol with a small amount of client storage'' achieving $O(\log N)$ bandwidth for blocks of
size $\Omega(\log^2 N)$ bits.

Classical ORAM assumes a client with plaintext visibility talking to an untrusted server. In
secure computation there is no such split: the computing parties are both client and server, and
no party may learn the access pattern. This is \emph{distributed} ORAM, introduced for two-party
computation by Lu and Ostrovsky~\cite{lu2013} as an alternative to unrolling a RAM program into a
circuit.

Floram~\cite{floram2017} made DORAM practical by building it from FSS. The authors report
improving prior constructions' ``access time by a factor of up to ten, their memory overhead by a
factor of one hundred or more, and their initialization time by a factor of thousands'', and
instantiating ORAMs holding $2^{34}$ bytes. They are explicit about the trade they make: deriving
the construction from FSS ``significantly reduces the amount of secure computation required to
implement an ORAM access, albeit at the cost of $O(n)$ efficient local memory
operations''~\cite{floram2017}. That trade, accepting linear local computation to buy sublinear
communication, recurs throughout Section~\ref{sec:retrieval}.

\duoram{} follows Floram's DPF-based approach and reports reducing communication further: for
memories with $n$ addressable locations it performs ``a sequence of $m$ arbitrarily interleaved
reads and writes'' using $O(m \lg n)$ words against Floram's $O(m\sqrt{n})$, with most of the work
moved into ``a data-independent preprocessing phase, leaving just $O(m)$ words of online
communication cost for the sequence --- i.e., a constant online communication cost per memory
access''. Other points in the design space include two-server constructions achieving ``sublinear
server computation and constant rounds''~\cite{hamlin2021}, and three-party protocols matching
client-server Path ORAM in round complexity~\cite{jarecki2018,bunn2020}.

Oblivious \emph{data structures} build on this: Wang et al.~\cite{wang2014} apply pointer-based
and locality-based techniques to ``maps, sets, priority-queues, stacks, deques'', and \prac{}
contributes binary search, heap and AVL-tree protocols over \duoramplain{}.

\subsection{Fixed-point arithmetic, and why it dominates the cost}
\label{sec:prim:fixedpoint}

The primitives above compute over rings and finite fields, while recommendation computes over the reals.
Bridging that gap with fixed-point arithmetic is where much of the practical cost of these systems
is incurred, which is why it deserves its own subsection rather than a footnote.

A real value $v$ is represented as $\lfloor v \cdot 2^{\fracbits} \rfloor \in \ring$. Addition is
then free, but the product of two $\fracbits$-scaled values is $2\fracbits$-scaled and must be
\emph{truncated} back. Repeated multiplication also requires \emph{normalisation} to stop values
growing without bound. Both are non-linear operations, and both are expensive under secret
sharing.

Consequently a recurring theme in Section~\ref{sec:training} is that the linear algebra is cheap
and the non-linear steps are not. SecureML~\cite{secureml2017} contributes ``new techniques to
support secure arithmetic operations on shared decimal numbers'' and ``MPC-friendly alternatives
to non-linear functions such as sigmoid and softmax''. Escudero et al.~\cite{escudero2020}
introduce \emph{extended doubly-authenticated bits}, ``shared integers in the arithmetic domain
whose bit decomposition is shared in the binary domain'', used to ``considerably increase the
efficiency of non-linear operations such as truncation, secure comparison and bit-decomposition''.
General-purpose frameworks make the same division visible: ABY~\cite{aby2015} and its
successor~\cite{aby22021} are \emph{mixed-protocol} systems precisely because different operations
are cheapest under different sharings, and MP-SPDZ~\cite{mpspdz2020} implements ``34 MPC protocol
variants'' behind one interface to make such comparisons possible.
 + %% ==========================================================================
%%  SECTION 2 --- BACKGROUND
%%  Every claim traces to report/notes/evidence.md.
%% ==========================================================================
\section{Background}
\label{sec:background}

A recommender system is a data-mining pipeline whose input is, by construction, a record of what
people consume. This section fixes the vocabulary the rest of the survey uses, and introduces the
device around which it is organised: the recommendation \emph{pipeline}, and the observation that
different bodies of work protect different stages of it.

\subsection{Collaborative filtering and matrix factorisation}
\label{sec:background:cf}

Collaborative filtering predicts a user's preferences from the preferences of other users.
Approaches divide broadly into neighbourhood methods, which score an item by its similarity to
items the user already liked, and latent-factor models, which embed users and items in a shared
low-dimensional space. Matrix factorisation is the dominant latent-factor
technique~\cite{koren2009,hu2008}. Given a partially observed user--item matrix
$\ratings \in \R^{\users \times \items}$, it seeks
%
\[
  \ratings \;\approx\; \userembed \cdot \itemembed,
  \qquad
  \userembed \in \R^{\users \times \dime},
  \quad
  \itemembed \in \R^{\dime \times \items},
  \quad
  \dime \ll \min(\users, \items),
\]
%
so that each user and each item is described by a $\dime$-dimensional embedding, and the predicted
affinity of user $i$ for item $j$ is the inner product of their embeddings.

Neural approaches to collaborative filtering exist and report improvements: Liang et
al.~\cite{liang2018} extend variational autoencoders to implicit feedback, describing a
``non-linear probabilistic model'' that goes ``beyond the limited modeling capacity of linear
factor models which still largely dominate collaborative filtering research'', and report
outperforming several baselines including two neural ones. That claim should be read alongside
Ferrari Dacrema et al.~\cite{dacrema2019}, whose systematic analysis of neural recommendation
proposals was motivated by ``problems in today's research practice \dots\ in terms of the
reproducibility of the results or the choice of the baselines''.

This tension matters for the present survey in a specific way. The cryptographic literature
surveyed in Section~\ref{sec:training} targets matrix factorisation almost exclusively, and does
not attempt to run deep recommenders under secure computation. Given that matrix factorisation
remains a competitive baseline, and that every additional non-linearity is expensive to evaluate
securely (Section~\ref{sec:prim:fixedpoint}), that focus is a reasonable engineering choice rather
than a gap.

\subsection{The recommendation pipeline}
\label{sec:background:pipeline}

It is useful to separate a deployed recommender into four stages:

\begin{enumerate}
  \item \textbf{Collection.} Ratings, clicks or views leave the user and reach the service.
  \item \textbf{Training.} A model --- here, the factorisation
        $\ratings \approx \userembed \cdot \itemembed$ --- is fitted to the collected data.
  \item \textbf{Serving.} For a given user, scores are computed over the catalogue and a ranking
        is produced.
  \item \textbf{Delivery.} The user fetches the recommended item itself: the film, the article,
        the product page.
\end{enumerate}

These stages are separable, and different lines of work protect different subsets of them. That
observation organises this survey, and Section~\ref{sec:taxonomy} makes it one of the four axes of
the taxonomy. It also carries an immediate consequence: \emph{a pipeline is only as private as its
weakest stage}. Training a model under secure computation achieves little if the subsequent fetch
of the recommended item is performed in the clear, because the fetch reveals the recommendation
that the training protocol was protecting.

\begin{figure}[t]
\centering
\begin{tikzpicture}[
    font=\small,
    stage/.style={draw, rounded corners=2pt, minimum width=21mm, minimum height=8mm,
                  align=center, inner sep=3pt},
    leak/.style={font=\scriptsize\itshape, align=center, text width=25mm},
    >={Stealth[length=5pt]}
  ]
  \node[stage] (c) at (0,0)    {Collection};
  \node[stage] (t) at (3.1,0)  {Training};
  \node[stage] (s) at (6.2,0)  {Serving};
  \node[stage] (d) at (9.3,0)  {Delivery};

  \draw[->] (c) -- (t);
  \draw[->] (t) -- (s);
  \draw[->] (s) -- (d);

  \node[leak, below=5mm of c] {ratings, and\\ who submitted them};
  \node[leak, below=5mm of t] {the training set;\\ intermediate model state};
  \node[leak, below=5mm of s] {which user asked;\\ what scored highly};
  \node[leak, below=5mm of d] {which item was\\ actually fetched};

  \draw[decorate, decoration={brace, mirror, amplitude=4pt}]
        ($(c.south west)+(-0.1,-1.55)$) -- ($(d.south east)+(0.1,-1.55)$)
        node[midway, below=5pt, font=\scriptsize] {what an adversary observes at each stage};
\end{tikzpicture}
\caption{The recommendation pipeline. Approaches surveyed here protect different subsets of these
         four stages; a pipeline is only as private as its weakest stage.}
\label{fig:pipeline}
\end{figure}

\subsection{What the adversary observes}
\label{sec:background:adversary}

Figure~\ref{fig:pipeline} annotates each stage with what is exposed to a service that runs it
without protection. Two results establish that this exposure is a demonstrated risk rather than a
hypothetical one.

Narayanan and Shmatikov~\cite{narayanan2008} present ``a new class of statistical
de-anonymization attacks against high-dimensional micro-data, such as individual preferences,
recommendations, transaction records''. Applying it to the Netflix Prize dataset --- ``anonymous
movie ratings of 500,000 subscribers'' --- with the Internet Movie Database as background
knowledge, they report having ``successfully identified the Netflix records of known users,
uncovering their apparent political preferences and other potentially sensitive information''.
They further note that the attack is ``robust to perturbation in the data'' and tolerates ``some
mistakes in the adversary's background knowledge''. Ratings data, in other words, is identifying
and not merely sensitive, and releasing it in a nominally anonymised form is not sufficient.

Calandrino et al.~\cite{calandrino2011} go further and consider the privacy risks of collaborative
filtering itself.\footnote{We were unable to obtain this paper: it is paywalled and no open
version was found. We therefore cite it only for the claim its title makes, and do not
characterise its method or results.} The distinction is important for
Section~\ref{sec:open:functionality}: an attack on the \emph{released dataset} is addressed by
protecting the pipeline, whereas an attack on the recommender's \emph{intended output} is not.

\subsection{Utility, and why it constrains the cryptography}
\label{sec:background:utility}

Recommendation quality is measured by ranking metrics computed over a held-out set, most commonly
normalised discounted cumulative gain (nDCG@$\topk$) and recall at $\topk$. This matters here
because a private recommender that recommends badly has solved nothing, and because the
constraint propagates into protocol design: fixed-point arithmetic, approximation of non-linear
functions, and reduced-precision representations must not perturb the ranking. Reporting quality
alongside cost is accordingly the norm in this literature, and
Section~\ref{sec:eval} returns to it.
.
%%  final.tex's original header said to reuse the survey's sections wholesale.
%%  They cannot be reused as they stand: they refer to themselves as "this
%%  survey" throughout and cross-reference sections that exist only in the
%%  survey (sec:training, sec:retrieval, sec:taxonomy, sec:open:*), which would
%%  print as "??" here. Editing them is not an option either, because they are
%%  the source of the submitted Milestone 1 PDF. So this file restates what the
%%  final report needs, with the same citations and the same notation.
%% ==========================================================================
\section{Background}
\label{sec:background}

This section fixes the primitives the rest of the report is built from. It is
deliberate about \emph{costs}, and in particular about rounds, because on a
wide-area link the round count decides the running time of every protocol here.

\subsection{Matrix factorisation by power iteration}
\label{sec:background:mf}

Collaborative filtering predicts a user's preferences from the preferences of
other users, and matrix factorisation is its dominant latent-factor
technique~\cite{koren2009,hu2008}. Given the rating matrix $\ratings \in
\R^{\users \times \items}$, it seeks $\ratings \approx \userembed \itemembed$
with $\userembed \in \R^{\users \times \dime}$, $\itemembed \in \R^{\dime \times
\items}$ and $\dime \ll \min(\users, \items)$. The predicted affinity of user $i$
for item $j$ is the inner product of their embeddings, and a user's
recommendations are the $\topk$ highest-scoring items the user has not already
rated.

\nudgeplain{} computes the factorisation by power iteration rather than by
gradient descent. For each of $\dime$ components it starts from a random vector
$\mathbf{v}$, and $\iters$ times it multiplies by $\ratings^{\top}\ratings$,
projects out the components already found, and normalises. The converged
$\mathbf{v}$ becomes a row of $\itemembed$, and finally $\userembed =
\ratings\itemembed^{\top}$. Power iteration on $\ratings^{\top}\ratings$
converges to the top-$\dime$ right singular subspace of $\ratings$, at a rate set
by the gap between consecutive singular values. What makes this formulation
attractive under secret sharing is its shape: almost all of the work is
matrix--vector products, with one normalisation per step.

\subsection{Replicated secret sharing}
\label{sec:background:rss}

In 2-out-of-3 replicated sharing over the ring $\ring = \Ztwo{\bits}$, a value
$x$ is written as $x = x_0 + x_1 + x_2$, and party $\party{i}$ holds the pair
$(x_i, x_{i+1})$ with indices taken modulo three~\cite{araki2016}. Any two parties
can reconstruct $x$, while a single party holds two uniformly random summands and
learns nothing. We write $\shr{x}$ for a sharing of $x$. Addition, and
multiplication by a public constant, are local.

Multiplication of two shared values is where the structure pays off. Because
each party holds two of the three summands of each factor, it can compute three
of the nine cross terms of $xy$ on its own, and between them the three parties
cover all nine. Each party therefore produces an additive share $z_i$ of $xy$
locally. To return to replicated form, each party masks $z_i$ with a share of
zero and sends it to one neighbour: one round and one ring element per party per
product. The zero shares cost no communication in the \emph{PRF model}, where
each pair of parties holds a shared pseudorandom-function key and the three
parties derive $\alpha_i = F(k_i, \mathit{ctr}) - F(k_{i-1}, \mathit{ctr})$,
which sum to zero~\cite{araki2016}.

The consequence \nudgeplain{} builds on is its Theorem~4.2. An inner product is a
sum of products, and the sum can be taken before resharing, so a whole inner
product costs one ring element per party, not one per term. A matrix--vector
product is one inner product per output entry. Its communication is therefore
proportional to the length of the \emph{output} vector, whatever the size of the
matrix.

\subsection{Fixed-point arithmetic and truncation}
\label{sec:background:fixed}

Recommendation works over the reals and the sharing works over a ring, so a real
$v$ is encoded as the integer $\lfloor v \cdot 2^{\fracbits} \rceil \in \ring$
with $\fracbits$ fractional bits. Addition then works unchanged, but the product
of two values at scale $\fracbits$ is at scale $2\fracbits$ and must be
\emph{truncated} back, or the scale grows until the ring overflows. Truncating
shares naively, by having each party shift its own share, fails: the shares are
uniform over the ring, so their sum wraps with high probability, and a wrap turns
into an error of $2^{\bits - \fracbits}$ in the result. Correct truncation needs
an interactive protocol. It is the most frequent non-linear step in fixed-point
secure computation, and prior work treats it as a first-class
cost~\cite{secureml2017,escudero2020}. Section~\ref{sec:eval:stages} shows that
it dominates ours.

\subsection{Function secret sharing}
\label{sec:background:fss}

Function secret sharing splits a \emph{function} rather than a
value~\cite{bgi2015,bgi2016}. A two-party scheme turns $f\colon \{0,1\}^{D} \to
\ring$ into keys $k_0, k_1$ such that each key alone hides $f$ and the two
evaluations recombine to $f(x)$ at every $x$. A \emph{distributed point function}
(DPF) is the case where $f = f_{\alpha,\beta}$ is $\beta$ at $x = \alpha$ and zero
elsewhere~\cite{gilboa2014}. The construction of Boyle, Gilboa and Ishai is a
binary tree of pseudorandom seeds, and each key holds one seed and one correction
word per level~\cite{bgi2016}. Keys are therefore \emph{logarithmic} in the domain
size $2^{D}$, while evaluating the whole domain costs linear work, a trade made
concrete in Section~\ref{sec:impl:dpf}.

A \emph{distributed comparison function} (DCF) shares the step function
$f_{<\alpha}(x) = \beta \cdot 1[x < \alpha]$ in the same way~\cite{boyle2021}. A
DCF key hides its threshold and is evaluated at a public point. Secure
computation needs the opposite, a secret input compared against a public
threshold, and Boyle et al.\ obtain it by masking. A dealer draws a random $r$
and builds keys for the threshold shifted by $r$. The parties open $c = x + r$,
which reveals nothing about $x$ if $r$ is drawn from a wide enough range, and
evaluate their keys at $c$. The evaluation yields shares of the comparison on
$x$~\cite{boyle2021}. Every threshold evaluated at the same opening costs no
further round.

\subsection{Two-server private information retrieval}
\label{sec:background:pir}

Private information retrieval lets a client fetch record $\alpha$ of a database
$D$ without the server learning $\alpha$~\cite{chor1995}. With two servers that
do not collude, a DPF gives a simple scheme~\cite{bgi2015,hafiz2019}. The client
generates keys for $f_{\alpha,1}$ and sends one to each server. Server $c$
expands its key over the whole domain and answers with
%
\[
  \mathrm{ans}_c \;=\; \sum_{j} \EvalFull(k_c)[j] \cdot D[j],
\]
%
and because the two expansions recombine to the indicator of $\alpha$, the two
answers recombine to $D[\alpha]$. The upload is one logarithmic-size key per
server. The price is on the server side: each answer touches every record, and
it must, because a server that skipped records would learn which ones the client
could not have wanted.

\subsection{Measuring recommendation quality}
\label{sec:background:ndcg}

We report normalised discounted cumulative gain at 20, nDCG@20, over the held-out
ratings of MovieLens-100K's predefined \texttt{u1} split~\cite{movielens2016},
with an item counted as relevant when its held-out rating is at least four and
with items already rated in training excluded from the ranking. The split has 943
users and 1{,}682 items. Our metric was cross-checked against an independent
implementation on two hundred random instances. Because the relevance convention
differs, our absolute values are a within-project baseline and are not comparable
to the numbers \nudgeplain{} reports on Netflix data.
